\documentclass[11pt]{article}
\usepackage[T1]{fontenc}
\usepackage{lmodern,amsmath,amssymb,amsthm,microtype}
\usepackage[margin=0.95in]{geometry}
\usepackage[colorlinks=true,allcolors=blue]{hyperref}
\newtheorem{theorem}{Theorem}
\newtheorem{lemma}{Lemma}
\newcommand{\CC}{\mathbb C}
\newcommand{\RR}{\mathbb R}
\newcommand{\area}{\mathcal A}
\title{A flag-curve proof of an original-side pedal-area invariant}
\author{Alec Kriebel\\
\small Independent researcher\\
\small \href{https://orcid.org/0009-0001-9320-500X}{ORCID: 0009-0001-9320-500X}}
\date{16 October 2026}
\begin{document}
\maketitle
\begin{abstract}
For a proper confocal elliptical billiard caustic with least period divisible
by four, the signed trajectory area times its signed pedal area is constant
along the family for every fixed pedal point. We present an incidence-curve
proof that includes primitive stars and zero or negative pedal areas.
The two factors have disjoint possible simple poles; reversal makes each
vanish at the other's poles. Compactness finishes the proof.
The exact theorem also appears in Ferudun's contemporary note, which we
credit. This is an alternative presentation using established Poncelet
methods, with no claim of theorem or method novelty, first disclosure,
or certified independent discovery. The least-period qualification is
essential, as an exact repeated-orbit control illustrates.
\end{abstract}

\section{Statement, conventions and prior work}
Fix the conics
\[
 E:\frac{x^2}{a^2}+\frac{y^2}{b^2}=1,\qquad
 C:\frac{x^2}{\alpha^2}+\frac{y^2}{\beta^2}=1,\qquad
 \alpha^2=a^2-\lambda,\quad\beta^2=b^2-\lambda,
 \quad a>b>0,\quad 0<\lambda<b^2 .
\]
Suppose the physical billiard trajectories with caustic \(C\) have
\emph{least} period \(N=2m\), with \(m\ge2\) even.
Write \(P_j\) for successive boundary vertices and \(L_j=P_jP_{j+1}\)
for their supporting lines, with indices modulo \(N\).
For each fixed finite \(M\in\RR^2\), let \(Q_j(M)\) be its perpendicular
foot on \(L_j\), and set
\[
 \area(Z)=\frac12\sum_{j=0}^{N-1}\det(Z_j,Z_{j+1}),
 \qquad A=\area(P),\quad B_M=\area(Q(M)).
\]
The order is the actual traversal order, including for a star.
These are signed shoelace areas. Feet are taken on complete lines,
with no clipping to side segments. The point \(M\) remains fixed as
the orbit phase varies.

\begin{theorem}\label{main}
For the fixed conic pair and each fixed \(M\), \(A B_M\) is independent
of phase. Every admitted coprime winding and both traversal orientations
are included. The conclusion allows zero or negative \(B_M\).
\end{theorem}

This is the least-period interpretation of \(k_{203,a}\) in
Reznik--Garcia--Koiller \cite[Section~3.3, Table~3, p.~346]{rgk};
the same row is in their arXiv edition \cite[Table~3, p.~6]{rgkarxiv}.
Their area convention is signed \cite[Eq.(1)]{rgk}. The table prints
\(N\equiv0\pmod4\) and all \(M\), but does not literally print
\emph{least period}. We state that qualification explicitly.
An unrestricted repeated-list interpretation is false
(Section~\ref{controls}). Historical question marks in the proof column
do not establish present-day openness.

Ferudun \cite[Theorem A, p.~2]{ferudun} states the same theorem, with
the stronger allowance of complex \(M\), the four-period formula and
signed star examples. His Section~4 proves the product by complementary
poles and reversal zeros in Jacobi coordinates. The proof below expresses
that strategy on the flag curve with explicit branch charts.
The original submitted proof preserved in our supplement uses a
pedal/contact-area trace comparison and the credited even area-product
theorem of Ch\'avez-Caliz \cite{chavez}. Different organizations do not
establish theorem or method novelty.

The radial pedal-area principle is classical; see Steiner's treatment
\cite[Sections XIV--XVI, reprint pp.~18--22]{steiner}.
The flag curve, involutions and meromorphic-invariant methods are
established Poncelet tools \cite[Section~3]{schwartz},
\cite[Sections~4--5]{ast}, \cite[Section~2 and p.~104]{chavez}.
The four-period case also follows from Bialy--Tabachnikov's distance-product
invariant \cite[Theorem~4.3(2)]{bt}, as explained below.
We use the flag construction in \cite{schwartz}, without relying on its
centroid conclusion whose proof gap is discussed in \cite[p.~92]{chavez}.

\paragraph{Dated records.}
The currently public original GitHub PR203 record gives server creation
time \(2026\text{-}10\text{-}01\), 05:24:27 UTC; its original head
\texttt{8edca9d773c085adb4ed0611357a7a2957e3395d}
contains the preserved proof \cite{pr203}.
The original Zenodo record for \cite{ferudun} was created on
\(2026\text{-}10\text{-}15\), 16:11:42.588744 UTC.
These are verified object-record dates, documented in the supplement.
Historical repository visibility at every earlier instant, discovery
times, earliest disclosure and independence are not certified.
We describe separate contemporary records and claim no firstness.

\section{Flags, translations and the central half-period}
The complex projective incidence curve is
\[
 X=\{(P,L):P\in E,\ L\text{ tangent to }C,\ P\in L\}.
\]
Let \(\pi\) forget \(L\), and \(\rho\) forget \(P\).
Let \(\sigma\) exchange the two endpoints of \(L\) on \(E\), and
\(\tau\) exchange the two caustic tangents through \(P\).
Composition is right to left; put \(T=\tau\sigma\).

On a physical orbit this is exactly the flag step. The reflected chord can
have the incoming chord's supporting line only at normal incidence:
a tangent to \(E\) cannot touch the strictly interior \(C\).
At \(P=(a\cos\theta,b\sin\theta)\), the normal line has coefficients
\[
 [u:v:w]=[-\sin\theta/b:\cos\theta/a:
                 c\sin\theta\cos\theta/(ab)],\qquad c=a^2-b^2 .
\]
Substitution in the dual-caustic equation gives
\(\lambda=a^2\sin^2\theta+b^2\cos^2\theta\ge b^2\), excluded.
Thus reflection selects the other caustic tangent.
The real-component description below also shows that vertex and flag
least periods agree.

We record the geometry and local charts needed for the proof.
Put \(c=a^2-b^2>0\). In line coordinates \(L=[u:v:w]\), the dual caustic is
\[
 C^*:\alpha^2u^2+\beta^2v^2-w^2=0,\qquad
 D=u^2+v^2,\quad H=a^2u^2+b^2v^2.
\]
Its endpoint double cover has the local homogeneous formula
\begin{equation}\label{cover}
 \eta^2=a^2b^2\lambda D,\qquad
 P=[-a^2uw-v\eta:-b^2vw+u\eta:H].
\end{equation}
Indeed \(H-w^2=\lambda D\); substitution verifies incidence and
\(x^2/a^2+y^2/b^2=z^2\). The formula extends by incidence at charts
where the displayed triple cancels.

The zeros of \(D\) on \(C^*\) are exactly
\begin{equation}\label{branch}
 L_{\epsilon,\delta}=[1:\epsilon i:\delta\sqrt c],
 \qquad \epsilon,\delta\in\{1,-1\}.
\end{equation}
They are the common tangents of \(E,C\). They are simple zeros:
on \(u=1\), \(w\ne0\), \(v\) is a coordinate and \(D'=2v\ne0\).
Thus \(\rho\) branches simply at these four values, and has no other
branch values. The smooth double cover of the conic is connected:
a disconnected degree-two cover would consist of degree-one covers
and could not branch. Riemann--Hurwitz gives genus one.
This also verifies smoothness of the incidence curve through the
charts omitted by the displayed endpoint formula.

The other projection \(\pi:X\to E\) branches precisely at the four
finite transverse intersections
\[
 x^2=\frac{a^2\alpha^2}{c},\qquad
 y^2=-\frac{b^2\beta^2}{c},\qquad z=1.
\]
All coordinates are nonzero and the gradients are not proportional,
since \(a^2/\alpha^2\ne b^2/\beta^2\).
In particular \(\pi\) is unramified over both infinities of \(E\).

Each branched deck involution acts by \(-1\) on a holomorphic
differential of \(X\). Hence \(T\) preserves that differential and is
a translation under genus-one uniformization. A real orbit of least
period \(N\) makes its order exactly \(N\) on \(X\).
Nonidentity powers of \(T\) have no fixed point. We have
\begin{equation}\label{reverse}
 \sigma T^j=T^{-j}\sigma,\qquad \tau T^j=T^{-j}\tau .
\end{equation}

For completeness, central inversion \(I\) is the half-period
\begin{equation}\label{central}
 I=T^m,\qquad N=2m .
\end{equation}
Here is a real-component proof that includes stars.
Normalize a real tangent as \(n(\phi)\cdot P=1\), where
\[
 n(\phi)=(\cos\phi/\alpha,\sin\phi/\beta),\quad
 q(\phi)=(\alpha\cos\phi,\beta\sin\phi),\quad
 d(\phi)=(-\sin\phi/\beta,\cos\phi/\alpha).
\]
Its endpoints are \(P=q+s d\).
Since \(q\) lies strictly inside \(E\), the quadratic equation for \(s\)
has distinct roots of opposite signs. These two continuous branches
give the two circles of \(X(\RR)\); \(\sigma\) exchanges them.
For fixed \(P\), the two tangent solutions of \(n(\phi)\cdot P=1\)
have derivative \(s/(\alpha\beta)\) of opposite signs, so \(\tau\)
also exchanges the circles. Each circle maps homeomorphically to
\(E(\RR)\), and \(T\) preserves it.

The map \(I\) sends \(\phi\) to \(\phi+\pi\) with \(s\) unchanged.
It commutes with \(T\), preserves either circle and its orientation,
and has order two. The map \(T\) preserves orientation, since it
preserves a nonvanishing real holomorphic differential there.
The finite group generated by \(T,I\) acts faithfully on that circle;
an automorphism identical on a real circle is identical on connected
\(X\). Averaging a metric over this finite group makes its
orientation-preserving elements rotations. There is just one
nonidentity rotation of order two. Thus \(I=T^m\) on the real circle,
and then on \(X\) by analytic continuation. No convex traversal
order has been assumed.

For \(x\in X\), define meromorphically
\(P_j(x)=\pi(T^jx)\), \(L_j(x)=\rho(T^jx)\), and their cyclic
areas. Cyclic reindexing makes both areas \(T\)-invariant.
All real phases have finite vertices and feet.

\section{Pole supports and their orders}
Let \(S_0\) be the four flags over the two infinities
\([a:\epsilon ib:0]\) of \(E\), and let
\(R_0=\operatorname{Fix}(\sigma)\), the four branch points in
\eqref{branch}. Write
\[
 S=\bigcup_{j=0}^{N-1}T^jS_0,\qquad
 R=\bigcup_{j=0}^{N-1}T^jR_0 .
\]
Only containment in these finite sets is asserted; a special \(M\)
may remove a prospective pole.

\begin{lemma}\label{polesA}
The only possible poles of \(A\) lie in \(S\), and each has order at
most one. If \(x\in S_0\), then
\(\sigma x=T^{m-1}x\).
\end{lemma}
\begin{proof}
Affine coordinates on \(E\) have simple poles at its two infinities.
Their pullbacks still have simple poles, since \(\pi\) is unramified
there. Two consecutive vertices cannot both be infinite. If they
are the distinct infinities, their line is \(z=0\), not tangent to \(C\).
If the two intersections coincide at an infinity, the line must be
the tangent to \(E\) there. Its coordinates
\([1/a:\epsilon i/b:0]\) do not satisfy the dual-caustic equation:
\[
 \frac{\alpha^2}{a^2}-\frac{\beta^2}{b^2}
 =\lambda(1/b^2-1/a^2)\ne0.
\]
Thus every determinant in \(2A=\sum_j\det(P_j,P_{j+1})\)
has at most one simple-pole factor. Multiple nonadjacent infinite
vertices do not increase the pole order of this sum.

For a flag over \([a:\epsilon ib:0]\), take \(v=1\).
Incidence and tangency give
\[
 u=-\epsilon ib/a,\qquad w^2=-\lambda c/a^2\ne0 .
\]
The involution \(\tau\) exchanges the two choices of \(w\).
Central inversion fixes the projective infinite point and exchanges
the same flags. Thus \(\tau x=Ix=T^mx\);
since \(\tau=T\sigma\), the claimed identity follows.
\end{proof}

The perpendicular-foot map, extended bilinearly to the complex curve,
is the homogeneous rational expression
\begin{equation}\label{foot}
 Q(M;[u:v:w])
 =M-\frac{uM_x+vM_y+w}{u^2+v^2}(u,v).
\end{equation}
It is regular where \(D\ne0\), including at \(w=0\).
The only possible poles on \(C^*\) are its four simple \(D\)-zeros.
There is no point of \(C^*\) with \(u=v=0\).
At each \(D\)-zero the pullback \eqref{cover} doubles its order.

\begin{lemma}\label{polesB}
The only possible poles of \(B_M\) lie in \(R\), and each has order
at most one, for every fixed finite real \(M\).
\end{lemma}
\begin{proof}
At \(x\in R_0\), choose the local coordinate \(t=\eta\), with \(u=1\).
Equations \eqref{cover} and the dual conic give convergent local branches
\[
 D=\frac{t^2}{a^2b^2\lambda},\quad
 v=\epsilon i\sqrt{1-\frac{t^2}{a^2b^2\lambda}},\quad
 w=\delta\sqrt{c+\frac{\beta^2t^2}{a^2b^2\lambda}}.
\]
Here \(\sigma(t)=-t\), while \(v,w\), and hence \(Q_0(t)\),
are even Laurent functions of order at most two.

Because \(\rho\sigma=\rho\) and \eqref{reverse} holds,
\begin{equation}\label{Qreverse}
 Q_j(\sigma y)=Q_{-j}(y).
\end{equation}
The indices with \(T^jx\in R_0\) are exactly \(j=0,m\):
indeed \(\sigma T^jx=T^{-j}x=T^jx\) precisely when \(N\mid2j\).
These two indices are nonadjacent, since \(m\ge2\).
Their foot germs are even; all other feet are regular.
By \eqref{Qreverse}, each difference
\(Q_{i+1}(t)-Q_{i-1}(t)\), for \(i=0,m\), is odd, regular
and vanishes at \(t=0\). The two incident edges combine as
\[
 \det(Q_{i-1},Q_i)+\det(Q_i,Q_{i+1})
 =\det\bigl(Q_i,Q_{i+1}-Q_{i-1}\bigr).
\]
An order-at-most-two even pole multiplied by an odd function vanishing
to order at least one has order at most one.
The two singular-vertex groups have disjoint edges.
When \(N=4\) they account for all four edges exactly once.
Vanishing leading coefficients for special \(M\) only lower the
pole order. Translating by powers of \(T\) treats all of \(R\).
\end{proof}

\section{Complementary zeros and the product}
\begin{lemma}\label{zeros}
The sets \(S,R\) are disjoint. The function \(A\) is regular and
vanishes on \(R\), and \(B_M\) is regular and vanishes on \(S\).
\end{lemma}
\begin{proof}
If an \(x\in S_0\) had \(T^jx\in R_0\), Lemma~\ref{polesA}
and \eqref{reverse} would give
\[
 T^jx=\sigma T^jx=T^{-j}\sigma x=T^{m-1-j}x,
 \qquad 2j\equiv m-1\pmod{2m}.
\]
The congruence is impossible: \(m\) is even. Translation invariance
of the two supports proves their disjointness.

For \(x\in R_0\), use \(\sigma x=x\), \(\pi\tau=\pi\) and
\(\pi\sigma=\pi T\). They give
\[
 P_j(x)=P_j(\sigma x)=P_{1-j}(x).
\]
All vertices are regular there, since \(R\cap S=\varnothing\).
The involution \(j\mapsto1-j\) reverses the cyclic polygon, so its
signed area is its own negative: \(A(x)=0\).

For \(x\in S_0\), \(\sigma x=T^{m-1}x\) and
\eqref{Qreverse} give
\[
 Q_j(x)=Q_{-j}(\sigma x)=Q_{m-1-j}(x).
\]
All feet are regular there. The index involution reverses the side
order while \(M\) is fixed, so \(B_M(x)=0\).
Cyclic invariance extends both zero statements to \(R,S\).
\end{proof}

\begin{proof}[Proof of Theorem~\ref{main}]
The two areas are meromorphic on connected compact \(X\).
At \(S\), \(A\) has at most a simple pole, removed by the regular
zero of \(B_M\). At \(R\) the roles are reversed.
Lemmas~\ref{polesA}--\ref{zeros} exhaust all possible poles.
Thus \(A B_M\) is holomorphic on \(X\), hence constant.
No division by an area or by a coefficient depending on \(M\) is used.

The argument applies on both real components. Reversing traversal
negates both signed areas and preserves the product.
The real-component proof of \eqref{central} included every primitive
star, so no winding restriction has been introduced.
\end{proof}

\section{Classical radial dependence and exact controls}\label{controls}
For any cyclic list of real side lines \(n_j\cdot X=h_j\), with
unit normals and independently arbitrary choices of their signs, set
\[
 d_j=\det(n_j,n_{j+1}),\qquad s_j=n_j\cdot n_{j+1}.
\]
Translation invariance of shoelace area and direct projection give
\[
 B_M=\frac12\sum_jd_j(h_j-n_j\cdot M)(h_{j+1}-n_{j+1}\cdot M)
     =B_0+\ell\cdot M+C_2|M|^2,
\]
where
\[
 \ell=-\frac12\sum_jd_j(h_jn_{j+1}+h_{j+1}n_j),\qquad
 C_2=\frac14\sum_j d_js_j
     =\frac18\sum_j\sin(2\Delta_j).
\]
For example, taking \(M=(r,0)\), the anisotropic quadratic terms
telescope by
\[
 \sin(\beta-\alpha)\cos\alpha\cos\beta
 =\tfrac14\{\sin2\beta-\sin2\alpha+\sin(2\beta-2\alpha)\}.
\]
This is the signed shoelace form of the classical pedal-area principle.
If \(C_2=0\), the linear term must still be retained for an arbitrary
crossing polygon; no finite curvature center is assumed.
Opposite side pairing in our even-period setting makes \(\ell=0\).
Theorem~\ref{main}, evaluated at \(M=0\) and a unit point, then gives
\[
 A B_M=K_0+K_2|M|^2
\]
with phase-independent real coefficients. This consequence is also
part of \cite[Theorem A]{ferudun}.

\paragraph{Four periods.}
For a parallelogram the alternating turns give \(C_2=0\).
The four-period family contains the orbit
\((a,0),(0,b),(-a,0),(0,-b)\):
through an axis endpoint reflection and central pairing force the
next vertex to the other axis. Tangency gives
\(\lambda=a^2b^2/(a^2+b^2)\). Its original area is \(2ab\);
its centered pedal rectangle has area
\(4a^3b^3/(a^2+b^2)^2\). Therefore, for the whole family,
\[
 A B_M=\frac{8a^4b^4}{(a^2+b^2)^2}=8\alpha^2\beta^2.
\]
This value is already in \cite[Theorem A]{ferudun}.
It is also an explicit earlier corollary of
\cite[Theorem~4.3(2)]{bt}: for a parallelogram
\(P=(U,V,-U,-V)\), writing \(D=\det(U,V)\),
\(s=|V-U|\), \(t=|V+U|\), gives
\[
 A B_M=\frac{8D^4}{s^2t^2}
       =8\prod_{j=0}^3\operatorname{dist}(0,L_j).
\]
The cited distance-product invariant completes that deduction.

\paragraph{Why listed length is insufficient.}
Take \(a=2,b=1,\lambda=4/9,M=0\).
Two physical primitive six-period phases are
\[
 \begin{split}
 H={}&((2,0),(4/3,\sqrt5/3),(-4/3,\sqrt5/3),\\
     &\hspace{12mm}(-2,0),(-4/3,-\sqrt5/3),(4/3,-\sqrt5/3)),\\
 J={}&((0,1),(-4\sqrt2/3,1/3),(-4\sqrt2/3,-1/3),\\
     &\hspace{12mm}(0,-1),(4\sqrt2/3,-1/3),(4\sqrt2/3,1/3)).
 \end{split}
\]
Direct projection gives \(A B_0=14000/729\) and \(14336/729\).
Repeating each list twice makes a 12-entry walk and multiplies
each area by two. The products become \(56000/729\) and \(57344/729\),
with nonzero difference \(448/243\).
Exact incidence, tangency, reflection signs and products are checked
in the supplement. Thus merely having a listed length divisible
by four is insufficient. Ferudun \cite[Section~1.4 and Section~6]{ferudun}
also discusses the repeated-orbit obstruction.

\paragraph{Signed zero and negative pedals.}
For \(a^2=32,b^2=27,\lambda=864/35\), put \(a=4\sqrt2,b=3\sqrt3\),
\(x=4a/5,y=3b/5\), and
\[
 W=((a,0),(x,y),(0,b),(-x,y),(-a,0),(-x,-y),(0,-b),(x,-y)).
\]
The orbit \(P_j=W_{3j\bmod8}\) is a physical primitive \(3/8\) star.
Exact projection gives
\[
 A=\frac{168\sqrt6}{5},\qquad
 B_M=\frac{96\sqrt6}{245}(13-|M|^2),\qquad
 A B_M=\frac{13824}{175}(13-|M|^2).
\]
This is a rescaling of the same star family in
\cite[Table~2, p.~14]{ferudun}, not a new example claim.
At \(|M|^2=13\) the signed pedal area is zero; beyond it the area
is negative. The zero does not require the feet to be collinear,
and unsigned filled-face area is a different quantity.

\section{Verification, scope and research disclosure}
The supplement preserves the original 18-file submitted proof and
its authorship/source records, source-first plans, independent
derivations, adversarial reports and exact file manifests.
One analytic audit verifies the submitted Jacobi/contact-trace proof;
another reconstructs the complete incidence-curve argument here.
A direct Euclidean family proves the general pedal quadratic and
the full four-period result, but explicitly leaves two general-period
scalar invariants unresolved. That partial route is not counted as
another universal proof.

Fresh exact controls include physical reflection signs, primitive
stars, zero/negative pedals, reversal, nonconfocal countercontrols and
the repeated-list obstruction. Separately, direct physical
high-precision calculations cover 39 caustics, including nearly
circular and elongated ellipses and every admitted coprime winding
in the selected periods. Those finite noninterval diagnostics do
not prove the theorem; the proof above does.

The original checkers use removable Python assertions.
Their fresh ordinary replays match the archived result bytes.
With Python optimization \texttt{-O}, all assertion enforcement
disappears while advertised counts survive; deliberately false
identities then pass. Fresh exception-based checks reject the
wrong coefficient both ordinarily and with optimization.
The supplement distinguishes advertised and enforced checks and
does not certify unknown historical runtime settings.

The result is restricted to the stated proper confocal elliptical
caustics and least-period convention. No hyperbolic or degenerate
caustic, filled-face-area or unrestricted repeated-list extension
is claimed. The bounded priority audit records positive overlaps,
access limits and the chronology qualification; it is not an
exhaustive novelty certificate.

AI tools were used extensively in proof development, mathematical
auditing, literature checks, computation and drafting.
This is an unrefereed preprint. The recorded AI adversarial reviews
do not constitute conventional human peer review.

\begin{thebibliography}{10}
\bibitem{rgk} D.~Reznik, R.~Garcia and J.~Koiller,
\emph{Fifty New Invariants of \(N\)-Periodics in the Elliptic Billiard},
Arnold Math.~J.~7 (2021), 341--355.
\href{https://doi.org/10.1007/s40598-021-00174-y}{doi:10.1007/s40598-021-00174-y}.
\bibitem{rgkarxiv} D.~Reznik, R.~Garcia and J.~Koiller,
\emph{Eighty New Invariants of \(N\)-Periodics in the Elliptic Billiard},
\href{https://arxiv.org/abs/2004.12497v11}{arXiv:2004.12497v11} (2020).
\bibitem{ferudun} A.~Ferudun,
\emph{Pedal Areas of Periodic Elliptic Billiard Trajectories with
Respect to an Arbitrary Point}, version~1.0, 15 October 2026,
unrefereed preprint.
\href{https://doi.org/10.5281/zenodo.23400541}{doi:10.5281/zenodo.23400541}.
\bibitem{steiner} J.~Steiner,
\emph{Ueber den Kr\"ummungs-Schwerpunkt ebener Curven} (1838),
Sections XIV--XVI, reprint pp.~18--22, ETH e-rara.
\href{https://doi.org/10.3931/e-rara-3816}{doi:10.3931/e-rara-3816}.
\bibitem{schwartz} R.~Schwartz and S.~Tabachnikov,
\emph{Centers of mass of Poncelet polygons, 200 years after},
Math.~Intelligencer 38 (2016), 29--34.
\href{https://arxiv.org/abs/1607.04766}{arXiv:1607.04766}.
\bibitem{ast} A.~Akopyan, R.~Schwartz and S.~Tabachnikov,
\emph{Billiards in ellipses revisited}, European J.~Math.~8 (2022),
1313--1327.
\href{https://arxiv.org/abs/2001.02934v2}{arXiv:2001.02934v2}.
\bibitem{chavez} A.~Ch\'avez-Caliz,
\emph{More About Areas and Centers of Poncelet Polygons},
Arnold Math.~J.~7 (2021), 91--106.
\href{https://doi.org/10.1007/s40598-020-00154-8}{doi:10.1007/s40598-020-00154-8}.
\bibitem{bt} M.~Bialy and S.~Tabachnikov,
\emph{Dan Reznik's identities and more}, European J.~Math.~8 (2022),
1341--1354.
\href{https://doi.org/10.1007/s40879-020-00428-7}{doi:10.1007/s40879-020-00428-7};
\href{https://arxiv.org/abs/2001.08469v2}{arXiv:2001.08469v2}.
\bibitem{pr203} A.~Kriebel, original submitted proof record,
\href{https://github.com/AlecKriebel/Math/pull/203}{GitHub PR203},
created 1 October 2026, original head identified in the text;
archived unchanged in the accompanying verification package.
\end{thebibliography}
\end{document}

